\documentclass{jsarticle}
\usepackage{ascmac,amsmath,amssymb,amsthm,bm,physics,arydshln}
\newtheorem{question}{問}
\begin{document}
\textbf{基礎解析学２B　レポート課題}
\begin{flushright}
6122111 三森誠真
\end{flushright}
\begin{question}
\end{question}
偽である. $\exists c\in\mathbb{R}$ s.t. $f(x)=c$ (定数関数) のとき, $f$は$C^2$級の関数であり, $f'(0) = 0$を満たすが, $x=0$は$f$の極値をとる点でも変曲点でもない.

($\Longrightarrow$ の反例)
$f(x) = 
\begin{cases}
 x^5 \sin(\frac{1}{x}) & (x \neq 0) \\
 0 & (x = 0)
\end{cases}$

($\Longleftarrow$ の反例)
$f(x) = x^3 - x$ など

\bigskip

\begin{question}
\end{question}
(1)
$f_1 (x,y) := x^2+y^4$

\bigskip

$\displaystyle\frac{\partial f}{\partial x}(x, y) = 2x$, \ 
$\displaystyle\frac{\partial f}{\partial y}(x, y) = 4y^3$ なので停留点は$(0,0)$.
$$
\begin{aligned}
& A_{f_1}(x, y)=\left(\begin{array}{cc}
\displaystyle\frac{\partial^2 f}{\partial x^2}(x, y) & \displaystyle\frac{\partial^2 f}{\partial x \partial^2 y}(x, y) \\
\displaystyle\frac{\partial^2 f}{\partial y \partial x}(x, y) & \displaystyle\frac{\partial^2 f}{\partial y^2}(x, y)
\end{array}\right)=\left(\begin{array}{cc}
2 & 2 x+4 y^3 \\
2 x+4 y^3 & 12 y^2
\end{array}\right) , \ H_{f_1}(0,0)=\left|\begin{array}{ll}
2 & 0 \\
0 & 0
\end{array}\right|=0
\end{aligned}
$$

となるが, $(x,y) \neq (0,0)$において$f_1(x,y)>f(0,0)$となるので$f_1$は$(0,0)$で極小値をとる.

\bigskip
\bigskip
\bigskip

(2)
$f_2 (x,y) := -x^2-y^4$

\bigskip

$\displaystyle\frac{\partial f}{\partial x}(x, y) = -2x$, \ 
$\displaystyle\frac{\partial f}{\partial y}(x, y) = -4y^3$ なので停留点は$(0,0)$.
$$
\begin{aligned}
& A_{f_2}(x, y)=\left(\begin{array}{cc}
\displaystyle\frac{\partial^2 f}{\partial x^2}(x, y) & \displaystyle\frac{\partial^2 f}{\partial x \partial^2 y}(x, y) \\
\displaystyle\frac{\partial^2 f}{\partial y \partial x}(x, y) & \displaystyle\frac{\partial^2 f}{\partial y^2}(x, y)
\end{array}\right)=\left(\begin{array}{cc}
-2 & -2 x-4 y^3 \\
-2 x-4 y^3 & -12 y^2
\end{array}\right) , \ H_{f_2}(0,0)=\left|\begin{array}{ll}
-2 & 0 \\
0 & 0
\end{array}\right|=0
\end{aligned}
$$

となるが, $(x,y) \neq (0,0)$において$f_2(x,y)<f(0,0)$となるので$f_2$は$(0,0)$で極大値をとる.

\bigskip
\bigskip
\bigskip

(3)
$f_3(x,y) := x^3+y^3$

\bigskip

$\displaystyle\frac{\partial f}{\partial x}(x, y) = 3x^2$, \ 
$\displaystyle\frac{\partial f}{\partial y}(x, y) = 3y^2$ なので停留点は$(0,0)$.
$$
\begin{aligned}
& A f_3(x, y)=\left(\begin{array}{cc}
6 x & 3 x^2+3 y^2 \\
3 x^2+3 y^2 & 6 y
\end{array}\right) , \ H_{f_3}(0,0)=\left|\begin{array}{ll}
0 & 0 \\
0 & 0
\end{array}\right|=0
\end{aligned}
$$

となるが, $x^3+y^3=(x+y)(x^2-xy+y^2)=(x+y)\left\{(x-\displaystyle\frac{y}{2})^2 + \frac{3}{4}y^2\right\}$ なので,

$x+y>0$のとき, $f_3(x,y)>f(0,0)$,

$x+y<0$のとき, $f_3(x,y)<f(0,0)$

となり$f_3$は$(0,0)$で極大値も極小値もとらない.

\end{document}
